\begin{figure}[H]
\centering
\begin{tikzpicture}[>=Stealth,scale=1.05]
  \draw[->] (-4,0) -- (6.5,0) node[right] {eje óptico};
  \filldraw[fill=cyan!15,draw=blue!60,thick]
    (0,-2.3) .. controls (-0.9,-1.2) and (-0.9,1.2) .. (0,2.3)
    .. controls (0.9,1.2) and (0.9,-1.2) .. cycle;
  \foreach \y in {-1.3,1.3}{
    \draw[black!55,thick] (-3.7,\y) -- (0,\y);
    \draw[blue,very thick] (0,\y) -- (3.0,0);
    \draw[red,very thick] (0,\y) -- (4.6,0);
  }
  \fill[blue] (3,0) circle (2pt);
  \fill[red] (4.6,0) circle (2pt);
  \draw[<->,blue] (0,-1.8) -- (3,-1.8) node[midway,below] {$f_{\text{azul}}$};
  \draw[<->,red] (0,-2.25) -- (4.6,-2.25) node[midway,below] {$f_{\text{rojo}}$};
  \draw[blue,opacity=.35,line width=3pt] (5.6,-0.65) -- (5.6,0.65);
  \draw[red,opacity=.35,line width=3pt] (5.85,-0.65) -- (5.85,0.65);
  \node[align=center] at (5.7,1.25) {bordes\\coloreados};
\end{tikzpicture}
\caption{Aberración cromática: cada longitud de onda posee una distancia focal distinta.}
\end{figure}